% !TeX program = xelatex
\documentclass[12pt]{nextart}

\UseTemplateSet{
  layout = en_doc,
  globalfonts = {cmu,shanggu},
  features = {tables,image,citations,hyperlinks}
}

\AddBibliographyResource{references.bib}

\title{English Article Template}
\author{Author Name}
\date{\today}

\begin{document}

\maketitle

\begin{abstract}
This template is intended as a practical IMPE LaTeX System article starter rather than a
minimal probe. It includes a title block, abstract, section structure, and a
few common document elements that can be rewritten directly into a real paper.
\end{abstract}

\section{Background}

Use this section to introduce the topic, research context, and source
materials. It is also a natural place to define the document scope and clarify
what kind of contribution the paper aims to make. A narrative citation can be
written with \textcite{sample2026}; a parenthetical citation can be written as
\parencite{sample2026}.

\section{Research Questions}

Typical framing questions might include:

\begin{enumerate}
  \item What is the textual or linguistic problem under discussion?
  \item How does the existing scholarship approach it?
  \item What new analysis or organization does this paper provide?
\end{enumerate}

\section{Sample Table}

\begin{table}[htbp]
  \centering
  \begin{tabular}{lll}
    \toprule
    Item & Description & Note \\
    \midrule
    Sources & Primary material & Replace as needed \\
    Method & Analysis and comparison & Expand later \\
    Result & Working observation & Refine in draft \\
    \bottomrule
  \end{tabular}
  \caption{Basic table example}
\end{table}

\section{Sample Figure}

\begin{figure}[htbp]
  \centering
  \rule{0.7\linewidth}{0.28\linewidth}
  \caption{Figure placeholder}
\end{figure}

\section{Conclusion}

The goal of this template is to provide a real writing starter with enough
structure to begin drafting immediately.

\PrintBibliography

\end{document}
